% Standalone problem document, generated from the adjacent README.md.
% Compile with XeLaTeX. Mathematical content is maintained in Markdown.
\documentclass[11pt,a4paper]{article}
\usepackage[fontsize=10.5pt]{scrextend}
\usepackage[margin=22mm,headheight=15pt,headsep=9mm,footskip=11mm]{geometry}
\usepackage{fontspec}
\setmainfont{texgyrepagella}[Extension=.otf,UprightFont=*-regular,BoldFont=*-bold,ItalicFont=*-italic,BoldItalicFont=*-bolditalic]
\setsansfont{texgyreheros}[Extension=.otf,UprightFont=*-regular,BoldFont=*-bold,ItalicFont=*-italic,BoldItalicFont=*-bolditalic]
\setmonofont{texgyrecursor}[Extension=.otf,UprightFont=*-regular,BoldFont=*-bold,ItalicFont=*-italic,BoldItalicFont=*-bolditalic,Scale=MatchLowercase]
\usepackage{amsmath,amssymb}
\usepackage{unicode-math}
\setmathfont{texgyrepagella-math.otf}
\usepackage{xcolor,microtype,enumitem,needspace,fancyhdr}
\usepackage{bookmark}
\definecolor{ink}{HTML}{17344B}
\definecolor{accent}{HTML}{197D80}
\definecolor{muted}{HTML}{596773}
\hypersetup{colorlinks=true,linkcolor=accent,urlcolor=accent,pdftitle={AC-05 - Strassen's
asymptotic rank conjecture for tight
tensors},pdfauthor={Open Problems in Numerical Linear Algebra}}
\urlstyle{same}
\setlength{\parindent}{0pt}
\setlength{\parskip}{5pt plus 1pt minus 1pt}
\linespread{1.04}
\setlength{\emergencystretch}{3em}
\widowpenalty=10000
\clubpenalty=10000
\displaywidowpenalty=10000
\allowdisplaybreaks[1]
\setlist{nosep,leftmargin=1.5em}
\providecommand{\tightlist}{\setlength{\itemsep}{0pt}\setlength{\parskip}{0pt}}
\providecommand{\pandocbounded}[1]{#1}
\pagestyle{fancy}
\fancyhf{}
\fancyhead[L]{\sffamily\footnotesize\color{muted}OPEN PROBLEMS IN NUMERICAL LINEAR ALGEBRA}
\fancyhead[R]{\sffamily\bfseries\footnotesize\color{accent}AC-05}
\fancyfoot[L]{\parbox[b]{0.9\textwidth}{\sffamily\fontsize{8}{10}\selectfont\color{muted}Editorial ratings; a bounded search cannot certify that a problem remains unsolved.\\Literature check: 2026-09-08}}
\fancyfoot[R]{\sffamily\footnotesize\color{muted}\thepage}
\renewcommand{\headrulewidth}{0.3pt}
\renewcommand{\headrule}{\hbox to\headwidth{\color{accent}\leaders\hrule height \headrulewidth\hfill}}
\makeatletter
\renewcommand{\section}{\Needspace{5\baselineskip}\@startsection{section}{1}{0pt}{12pt plus 2pt minus 2pt}{4pt}{\sffamily\bfseries\large\color{ink}}}
\renewcommand{\subsection}{\Needspace{4\baselineskip}\@startsection{subsection}{2}{0pt}{9pt plus 1pt minus 1pt}{3pt}{\sffamily\bfseries\normalsize\color{ink}}}
\makeatother
\setcounter{secnumdepth}{0}
\begin{document}
{\sffamily\small\color{muted}Arithmetic and complexity\par}
\vspace{3pt}
{\raggedright\hyphenpenalty=10000\exhyphenpenalty=10000\sffamily\bfseries\fontsize{21}{24}\selectfont\color{ink}Strassen's
asymptotic rank conjecture for tight tensors\par}
\vspace{7pt}
{\sffamily\small\textbf{Difficulty:} extreme\quad\textbf{Importance:} broadly
interesting\par}
\vspace{4pt}
{\color{accent}\hrule height 0.8pt}
\vspace{6pt}
\textbf{Topic:} complexity of structured tensor powers

\textbf{Status:} open; admitted after a bounded literature search found
no resolution.

\subsection{Context and notation}\label{context-and-notation}

Over \(\mathbb C\), the tensor rank \(R(T)\) is the minimum number of
pure tensors in an exact sum for \(T\).

\subsection{Problem statement}\label{problem-statement}

A tensor \(T\in(\mathbb C^d)^{\otimes3}\) is \emph{concise} if each of
its three matrix flattenings has rank \(d\). It is \emph{tight} if some
choice of bases admits injective functions
\(a,b,c:\{1,\ldots,d\}\to\mathbb Z\) with \(a(i)+b(j)+c(k)=0\) whenever
the coefficient \(T_{ijk}\) is nonzero. Is

\[\lim_{m\to\infty}R(T^{\otimes m})^{1/m}=d\]

for every positive integer \(d\) and every concise tight \(T\)? Tensor
powers use corresponding-factor grouping; rank is over \(\mathbb C\).

\subsection{Why it matters}\label{why-it-matters}

This is a structural conjecture about many bilinear computations, beyond
the single matrix-multiplication family.
\href{https://github.com/ajt60gaibb/OpenProblemsInNLA/blob/main/arithmetic-and-complexity/AC-04/README.md}{AC-04}
is a named special tensor with independent literature treatment; this
universal statement is not its equivalent reformulation.

\subsection{References and status}\label{references-and-status}

A. Björklund and P. Kaski,
\href{https://arxiv.org/abs/2310.11926}{\emph{The Asymptotic Rank
Conjecture and the Set Cover Conjecture are not Both True}}, Conjecture
3 and §2.3, states the conjecture and gives a conditional consequence,
not a refutation. K. Lee,
\href{https://arxiv.org/abs/2601.08119}{\emph{Asymptotic rank bounds: a
numerical census}} (2026), Conjecture 1, retains the tight/concise
formulation. Searches for ``asymptotic rank conjecture proved 2026''
found no resolution. Numerical evidence is not a proof.
\textbf{Admitted: no resolution located.}
\end{document}
